% ====================================================================
%  Backup slides (after the references; not part of the 45-min talk)
% ====================================================================

\begin{frame}[noframenumbering,plain]
\centering\vfill
{\usebeamercolor[fg]{structure}\Large\bfseries Backup slides}
\vfill
\end{frame}

% ---------------------------------------------------------- Background
\begin{frame}{Good ranking $\neq$ good probabilities}
\small
Three models, four patients (stars = true event times):
\begin{itemize}
  \item $M_1$: perfect ranking, poor calibration
  \item $M_2$: poor ranking, perfect calibration
  \item $M_3$: both
\end{itemize}
\begin{center}
\includegraphics[width=0.74\linewidth]{discrimiation_vs_calibration.png}
\end{center}
\vspace{-0.6em}
\figsource{Figure courtesy of Qi et al.~\cite{qi2025enhancing}}
\takeaway{A high C-index says nothing about whether the predicted times or probabilities are right.}
\end{frame}

\begin{frame}{Correcting for censoring: IPCW}
\small
A censored instance passes its weight on to those still at risk:
\[
  w_i = \frac{1}{\hat S_C(t_i)}, \qquad \hat S_C = \text{Kaplan--Meier of the censoring times}
\]
Uno's C-index \cite{uno2011c} and the Brier score / IBS \cite{graf1999assessment} reweight with $w_i$ -- assuming random censoring (RQ3).
\vspace{0.2em}
\begin{center}
\includegraphics[width=0.84\linewidth]{ipcw.pdf}
\end{center}
\vspace{-0.6em}
\figsource{Figure courtesy of Qi et al.~\cite{qi_current_2025}. Top: KM of $S_C(t)$. Bottom: weights passed on at each censoring.}
\end{frame}

% ---------------------------------------------------------------- RQ1
\begin{frame}{TMJ in JIA: coverage on average only}
\begin{columns}[c]
  \begin{column}{0.46\textwidth}
    \small
    1,035 children with JIA, 6,154 clinical examinations. A random forest flags examinations with likely TMJ involvement.
    \vspace{0.5em}
    \begin{itemize}
      \item \textbf{Early stage} (first 2 years): precision 0.86, sensitivity 0.70 -- an earlier clinical study reported 0.47.
      \vspace{0.3em}
      \item \textbf{Conformal prediction} \cite{vovk2005} adds a set of likely classes that should contain the true class 90\% of the time.
    \end{itemize}
  \end{column}
  \begin{column}{0.52\textwidth}
    \centering
    \includegraphics[width=0.88\linewidth]{plot_tmj_coverage.pdf}
    \figsource{Data from Christensen et al.~\cite{christensen_explainable_2024}}
  \end{column}
\end{columns}
\takeaway{The 90\% holds on average -- early-stage involvement is covered only 74\% of the time.}
\end{frame}

% ---------------------------------------------------------------- RQ2
\begin{frame}{Ablations: both components help}
\begin{columns}[c]
  \begin{column}{0.44\textwidth}
    \small
    \begin{itemize}
      \item \textbf{Joint training}: within-patient ordering improves -- 90.2 vs.\ 87.8 on Rotterdam.
      \item Independent training is slightly better on \emph{global} metrics.
      \item \textbf{Trajectory loss}: small gains in local CI on Rotterdam and EBMT -- small relative to run-to-run variability.
    \end{itemize}
  \end{column}
  \begin{column}{0.54\textwidth}
    \centering
    \includegraphics[width=\linewidth]{plot_mensa_ablation.pdf}
    \figsource{Data from Lillelund et al.~\cite{lillelund_mensa_2026}}
  \end{column}
\end{columns}
\end{frame}
